\section{Reward Augmentation}
\label{sec:reward_augmentation}

During training, the policy optimizer receives a total reward formed by augmenting the environment reward with a clipped intrinsic term:
\begin{equation}
    r_t
    = r_t^{\mathrm{ext}}
    + \eta_t\,\mathrm{clip}\!\left(r_t^{\mathrm{int}},-r_{\max},r_{\max}\right),
    \label{eq:reward_augmentation}
\end{equation}
where $\eta_t\ge 0$ controls the strength of intrinsic shaping and $r_{\max}>0$ bounds the magnitude of the intrinsic contribution.
The clipping operation prevents unusually large intrinsic values from dominating the advantage estimates used by the policy update.
For environment-terminal transitions, the intrinsic term is set to zero so that episode termination itself is not treated as an intrinsically rewarding event.

The augmented training objective can be written as
\begin{equation}
    J_{\mathrm{aug}}(\theta)
    = \mathbb{E}_{\tau\sim\pi_\theta}
    \left[\sum_{t=0}^{H-1}\gamma^t
    \left(r_t^{\mathrm{ext}}+\eta_t\,\bar{r}_t^{\mathrm{int}}\right)
    \right],
\end{equation}
where $\bar{r}_t^{\mathrm{int}}=\mathrm{clip}(r_t^{\mathrm{int}},-r_{\max},r_{\max})$.
Policy-gradient updates such as PPO optimize this augmented objective through advantage estimates computed from $r_t$ \cite{schulman-2017}.
Consequently, the intrinsic term affects the data distribution collected by the policy, the value targets used by the critic, and the direction of policy improvement.

The coefficient $\eta_t$ may be constant or scheduled over the course of training.
A time-varying coefficient is useful when intrinsic rewards are intended to encourage early exploration while allowing the extrinsic task objective to dominate later learning.
Let $p_t\in[0,1]$ denote normalized training progress.
A generic schedule can be represented as $\eta_t=\eta w(p_t)$, where $\eta$ is a base scale and $w(p_t)\in[0,1]$ is a nonincreasing weighting function.
The theoretical framework only requires $\eta_t$ to be nonnegative and bounded.

Reward augmentation introduces an explicit tradeoff.
If $\eta_t$ is too small, the intrinsic reward may not change exploration behavior in sparse-reward regions.
If $\eta_t$ is too large, the agent may optimize the auxiliary signal at the expense of the task reward.
The subsequent components are therefore designed so that $r_t^{\mathrm{int}}$ is not merely large in novel or unpredictable states, but large primarily when the transition indicates controllable feature-space change and local model improvement.